\documentclass[twoside, 10pt]{article}
\usepackage{geometry}
\geometry{bottom=4em,top=6em, inner=2.2cm, outer=4em,  headheight=\paperheight}
\usepackage[export]{adjustbox}
\usepackage{array}
\usepackage{amsmath}
\usepackage{amsfonts}
\usepackage{fancyhdr}
\usepackage{luacode}
\pagestyle{fancy}
\fancyhf{}
\lhead{Precalculus - BASE}
\chead{Quadratic v.s. Linear-Logarithmic}
\directlua{pageNum = 1}
\rhead{Practice, Page \directlua{tex.print(pageNum); pageNum = pageNum + 1; if pageNum > 3 then pageNum = 1 end}}
\usepackage{lastpage}
\usepackage{xcolor}
\usepackage{enumitem}
\usepackage{pifont}
\usepackage{graphicx}
\graphicspath{{../img}}
\usepackage{pgfplots}
\pgfplotsset{compat=1.18}
\usepackage{tabularx}
\usepackage{polynom}

\newcommand{\R}{\mathbb R}
\newcommand{\e}{{\rm e}}
\newcommand{\pobr}[1]{\left\langle#1\right\rangle}
\newcommand{\norm}[1]{\lVert #1 \rVert}
\newcommand{\abs}[1]{\lvert #1 \rvert}

\DeclareMathOperator{\xd}{d\!}
\DeclareMathOperator{\proj}{proj}

\title{}
\date{}

\newcommand{\template}{
\directlua{tex.print(utils.genVersion())}
\noindent
{\large
First Name \rule{6em}{.1pt}\hspace{\stretch{1}} Last Name \rule{6em}{.1pt}\hspace{\stretch{1}} Date \rule{1.5em}{.1pt} -- \rule{1.5em}{.1pt} -- \rule{1.5em}{.1pt}\hspace{\stretch{1}} Period \rule{2em}{.1pt}\hspace{\stretch{1}} Score \rule{2em}{.1pt}
}
\vspace{1em}
}

\begin{luacode*}
utils = require("../utils")
\end{luacode*}

\begin{document}
\template

Consider the following sequence of 10 numbers:
\[
42,\; 7,\; 56,\; 18,\; 3,\; 63,\; 25,\; 78,\; 12,\; 31
\]

Answer the questions below.

\begin{enumerate}[label=\arabic*.]
    \item Design an algorithm to sort the sequence from smallest to largest.  
    Clearly describe each step of your algorithm, and justify that the algorithm always terminates.
    \vspace{\stretch{1}}

    \item How many merge steps does your algorithm perform before the sequence is fully sorted?  
    Explain how you determined this number.
    \vspace{\stretch{1}}

    \item Suppose the length of the sequence is \(n\) instead of 10.  
    As \(n\) increases, how does the number of steps required by your algorithm change?  
    Describe the relationship using words and, if possible, mathematical notation.
    \vspace{\stretch{1}}
\end{enumerate}

\clearpage

Here is the same sequence of numbers:
\[
42,\; 7,\; 56,\; 18,\; 3,\; 63,\; 25,\; 78,\; 12,\; 31
\]

Notice that the sequence is \emph{not} ordered.  
To efficiently sort the sequence, we can apply the \emph{merge sort} algorithm.

\subsection*{Description of Merge Sort}

Merge sort works as follows:
\begin{enumerate}[label=(\roman*)]
    \item Divide the sequence into two smaller subsequences of approximately equal length.
    \item Recursively apply merge sort to each subsequence.
    \item Merge the two sorted subsequences into a single sorted sequence by repeatedly selecting the smaller leading element.
    \item Continue this process until the entire sequence is sorted.
\end{enumerate}

Answer the questions below.

\begin{enumerate}[label=\arabic*.]
    \item Apply the merge sort algorithm to the given sequence.  
    Show the result of each division and each merge step.
    \vspace{\stretch{1}}    
    \clearpage
    \item How many times must the sequence be divided until all subsequences contain exactly one element?
    \vspace{\stretch{1}}
    \item At a given level of division, how many total numbers are involved in the merging process?
    \vspace{\stretch{1}}
    \item Does this total number of elements involved in merging depend on the level?  
    Explain your reasoning.
    \vspace{\stretch{1}}
    \item Using your answers above, describe how the total amount of work required by merge sort changes as the length \(n\) of the sequence increases.
    \vspace{\stretch{1}}
\end{enumerate}

\end{document}
\end{document}